\subsection{Two further literature routes and their precise limits}

\paragraph{KT-STABLE-CONGRUENCE-ROUTE-AUDIT.}
O.~Randal-Williams, \emph{Stable cohomology of congruence
subgroups}, Compositio Math.~162 (2026), 1274--1317,
\url{https://doi.org/10.1017/S0010437X26103315},
Theorem~A, p.~1275, factors congruence-subgroup cohomology
as an exterior algebra tensor completed cohomology in
degrees below $p-1$. Equation~(7.3), p.~1296, identifies
the latter with the cohomology of the fibre of the
completion map in $K$-theory. The remaining factor is
therefore arithmetic input, not an independently computed
vanishing. Section~7.3.2, p.~1297, explicitly retains the
unknown finite $K_{4k}$ terms. In \S7.5, p.~1301, those
terms occur in the fibre's exact sequence; the subsequent
examples assume Vandiver and additional low-degree
Bernoulli units. Thus applying the displayed factorization
does not eliminate the target group. This is a
statement-and-formula scope check of Theorem~A and
\S\S7.2--7.5, not a rereading of the paper's full proof.

\paragraph{KT-RANKIN-COHEN-NONVANISHING-SCOPE.}
R.~Takloo-Bighash, \emph{Simultaneous nonvanishing of
quadratic twists via Rankin--Cohen brackets},
\url{https://arxiv.org/abs/2609.19649}, Theorem~1.1 and
Corollary~1.2 in \S1, keeps a discriminant and a desired
number of forms fixed while the weight grows. Its
conclusion is nonzero complex central values for enough
level-one forms. It does not assert unit values at a
fixed residue prime, nor the three Dirichlet values at
one varying negative discriminant needed in (FV.8).
Consequently this theorem does not discharge that
simultaneous modular nonvanishing condition. Only the
stated theorem, definitions and scope were checked;
no full proof check or counterexample to its theorem
is claimed.

Both attempts leave the all-weight arithmetic gate
unchanged and are excluded from progress-review movement.
